% Generated draft: requires mathematical and source review.
\documentclass[11pt]{article}
\usepackage[margin=25mm]{geometry}
\usepackage{fontspec}
\setmainfont{Latin Modern Roman}
\usepackage{amsmath,amssymb,mathtools,unicode-math}
\setmathfont{Latin Modern Math}
\usepackage{enumitem,needspace,xurl,hyperref}
\setlength{\parindent}{0pt}
\setlength{\parskip}{5pt}
\setlength{\emergencystretch}{3em}
\newcommand{\N}{\mathbb{N}}
\newcommand{\Z}{\mathbb{Z}}
\newcommand{\Q}{\mathbb{Q}}
\newcommand{\R}{\mathbb{R}}
\newcommand{\C}{\mathbb{C}}
\newcommand{\F}{\mathbb{F}}
\newcommand{\A}{\mathbb{A}}
\newcommand{\PP}{\mathbb{P}}
\newcommand{\E}{\mathbb{E}}
\newcommand{\eps}{\varepsilon}
\newcommand{\dd}{\,\mathrm{d}}
\begin{document}
\section*{UnsolvedMath ID: 701083}
\textbf{Source problem code:} catalogue entry
\subsection{Definitions and mathematical statement}
All groups considered are finite. A formation \(\mathfrak F\) is a nonempty class of groups closed under isomorphisms, homomorphic images, and subdirect products. It is \(S\)-closed if every subgroup of every member belongs to \(\mathfrak F\), and \(S_n\)-closed if every subnormal subgroup of every member belongs to \(\mathfrak F\). Here \(H\) is subnormal in \(G\) if there is a chain \(H=H_0\trianglelefteq H_1\trianglelefteq\cdots\trianglelefteq H_r=G\).

An \(\mathfrak F\)-subgroup means a subgroup belonging to \(\mathfrak F\). Write \(R_{\mathfrak F}(H,G)\) for the source's assertion that \(H\) is \(\mathfrak F\)-subnormal in \(G\). The supplied record does not define this relative subnormality notion or specify its chain convention. Consequently, this predicate is retained explicitly; its intended definition must be checked before this draft is treated as a self-contained formulation.

The source calls \(\mathfrak F\) \emph{superradical} if it is \(S_n\)-closed and satisfies
\[
\begin{gathered}
\text{for every finite group }G\text{ and all }A,B\leq G,\\
G=AB,\quad A,B\in\mathfrak F,\quad
R_{\mathfrak F}(A,G),\quad R_{\mathfrak F}(B,G)
\quad\Longrightarrow\quad G\in\mathfrak F.
\end{gathered}
\]
Here \(AB=\{ab:a\in A,\ b\in B\}\); the equality \(G=AB\) is a setwise product condition, with no normality assumption on either factor.

Using the usual chief-factor formulation, a formation is \emph{local} if there is a function \(f\) assigning a formation, or the empty class, to each prime such that
\[
\mathfrak F=
\left\{G:\ 
G/C_G(H/K)\in f(p)
\text{ for every chief factor }H/K\text{ of }G
\text{ and every prime }p\mid |H/K|
\right\}.
\]
A chief factor \(H/K\) has \(K,H\trianglelefteq G\), \(K<H\), and no normal subgroup of \(G\) strictly between \(K\) and \(H\). Its centralizer is
\[
C_G(H/K)=\{g\in G:ghg^{-1}h^{-1}\in K\text{ for every }h\in H\}.
\]
This formulation includes nonabelian chief factors; the supplied problem does not restrict its universe to soluble groups.

\textbf{Part (a), the remaining classification problem.} Determine all local formations \(\mathfrak F\) of finite groups that are superradical in the sense above. A complete answer should give necessary and sufficient conditions for a local formation to have this property, with the intended meaning of \(R_{\mathfrak F}\) fixed.

\textbf{Part (b), the archived historical assertion.} For a finite group \(G\), let \(R(G)\) be its largest soluble normal subgroup and let \(\Phi(R(G))\) be the intersection of the maximal subgroups of \(R(G)\), with \(\Phi(1)=1\). In the usual terminology, \(\mathfrak F\) is \emph{solubly saturated} if
\[
\text{for every finite group }G,\qquad
G/\Phi(R(G))\in\mathfrak F\ \Longrightarrow\ G\in\mathfrak F.
\]
Part (b) originally asked for a proof that every \(S\)-closed superradical formation is solubly saturated. The supplied Archive of Solved Problems entry instead reports that a counterexample was constructed by S. Yi and S. F. Kamornikov in 2016. Thus this is preserved as an archived assertion reported to have been refuted, not as a theorem to prove. The supplied registry labels the whole record \texttt{partially\_solved}, retaining part (a) as remaining work; these are supplied curation claims, not independently verified conclusions.
\subsection{Short English statement}
For Kourovka Notebook Problem 14.99, proposed by L. A. Shemetkov, part (a) asks for a classification of local formations of finite groups satisfying the stated superradical product condition. Part (b) is archived: the supplied source reports a counterexample to its claim that every subgroup-closed superradical formation is solubly saturated. The intended definition of \(\mathfrak F\)-subnormality still requires verification.
\subsection{Sources}
\begin{itemize}
\item Supplied local UnsolvedMath registry record: permanent ID \(701083\), exact source problem code \texttt{catalogue entry}, title \emph{Kourovka Notebook Problem 14.99}. The supplied record does not identify its exact file within the local registry.
\item L. A. Shemetkov, proposer of Problem 14.99, \emph{The Kourovka Notebook (Unsolved Problems in Group Theory)}, 14th issue (1999). The supplied background identifies an update in the 21st issue (2026), arXiv identifier \texttt{1401.0300v48}. The supplied collection description names E. I. Khukhro and V. D. Mazurov as editors. No external source was consulted for this draft.
\item Supplied Archive of Solved Problems attribution for part (b): S. Yi and S. F. Kamornikov, \emph{Siberian Math. J.}, volume 57, no. 2 (2016), pp. 260--264. The supplied entry states that a counterexample is constructed; it supplies no article title or link.
\item Supplied status review dated \texttt{2026-10-16}, with evidence label \texttt{published}: canonical status \texttt{partially\_solved}; part (b) is recorded as solved in the archive and the remaining part as retained in the main section. This assessment is supplied dated curation, not independently verified literature assessment.
\end{itemize}
\end{document}
